\section{Related Work}
\label{sec:related}

The right ancestor of \sys{} is not another consumer of a kernel packet
API; it is reactive flow processing. Open vSwitch classifies in
userspace and caches the result in a kernel flow table, and Pfaff et
al.\ are explicit about what that split costs when the cache
misses~\cite{pfaff2015ovs}. NVP realized distributed logical networks
on that substrate~\cite{koponen2014nvp}, and OVN's stateful ACLs use
\texttt{allow-related} backed by connection tracking, so a decision
taken on a connection's first packet is inherited by the rest of
it~\cite{ovnacl}. Userspace miss handling, cached verdicts, and
conntrack-backed stateful filtering are all prior art, and this paper
claims none of them. What it claims is narrower: evaluation of
Kubernetes NetworkPolicy semantics in userspace at connection
admission, without materializing the permitted-peer set at each
endpoint in advance, with local Pod metadata taken from the container
runtime rather than the API server. The novelty, such as it is, lies in
the placement and the object cached, not in the mechanism.

Andromeda and ONCache are the nearest treatments of what to cache and
when to invalidate it. Andromeda separates its flow-processing paths so
that provisioning, isolation, performance, and feature velocity can be
traded independently~\cite{dalton2018andromeda}; ONCache finds work
repeated across the layers of a container overlay, filtering included,
and caches the invariant part in eBPF~\cite{oncache2025}. \sys{} caches
one object, the per-connection verdict, and its invalidation story is
the strict-mode conntrack scan of Section~\ref{sec:design}. KUBEDIRECT
names API-mediated controller communication as a serverless scaling
bottleneck and builds state management for a more direct
path~\cite{kubedirect2026}; the same observation motivates \sys{}'s use
of NRI, which removes one asynchronous round trip for one class of
local facts and provides none of KUBEDIRECT's controller coordination.
Role-based microsegmentation infers roles and policies from observed
traffic~\cite{mani2025microsegmentation}; it produces the policy an
engine then has to enforce and says nothing about the cost of enforcing
it.

The point most likely to be lost in Section~\ref{sec:model} is that
eager per-peer expansion is one kernel-resident representation among
several, and not the one the other kernel-based enforcers cited here
use. Open
vSwitch's conjunctive match installs a rule with $n$ sources and $m$
destinations as $n+m$ flows sharing a \texttt{conj\_id} rather than
$n\times m$ flows~\cite{ovsfields}, and Antrea compiles NetworkPolicy
exactly this way, marking established connections in conntrack so that
only new flows walk the policy tables~\cite{antreapipeline}: the same
evaluate-once-per-connection property this paper attributes to NFQUEUE
plus conntrack, obtained without leaving the kernel. Calico's Felix
compiles selectors into ipsets referenced by a bounded number of
iptables or nftables rules, so per-endpoint state grows with rules
rather than peers~\cite{calicoarch}. nftables sets offer the same
factoring to anyone; Sutter's measurements~\cite{sutter2017nftables}
and kube-proxy's incremental nftables updates~\cite{winship2025nftables}
exist because lookup cost and update cost are separate questions. eBPF
maps are likewise keyed state that userspace updates
incrementally~\cite{linuxbpfmaps}; the $2I$ entries per endpoint of
Equation~\ref{eq:materialization} describe how Cilium chooses to key
its per-endpoint policy map by security
identity~\cite{ciliumidentities} under a \texttt{--bpf-policy-map-max}
that defaults to 16,384, tops out at 65,536, and can optionally lock
the endpoint down on overflow~\cite{ciliumagentref}. That is one
representation under one configuration. A capacity comparison that
means to say something about kernel enforcement as a class needs a
factored-kernel baseline, and this paper's does not have one.

Finally, a distribution tier between the datastore and the node agents
is established practice: Calico's Typha holds datastore state,
deduplicates events, and fans updates out so that each agent does not
carry its own load against the datastore~\cite{calicotypha}. IPTracker
has to be judged on its representation, caching, and recovery behavior;
the tier itself is not a contribution.

\section{Discussion and Future Work}
\label{sec:discussion}

\subsection{Compact Metadata Distribution with IPTracker}
\label{sec:iptracker}

Declining to run a distributed identity allocator does not make
distributed metadata go away. If $N$ agents each watch every Pod and
Pods change at rate $U$, the API server delivers on the order of $NU$
events; filtering, the watch cache, and event size move the constant of
that first-order count, not the product. IPTracker, introduced in PR
190~\cite{iptracker190}, moves the full Pod watch to $R$ distributors
and exposes compact IP-to-Pod records to the agents, so full Pod
subscriptions against the API server fall from $N$ to $R$. Other
subscriptions remain, and the fan-out moves rather than disappears: the
inspected client watches the entire compact-record prefix, a smaller
object delivered to the same number of subscribers rather than a
selective subscription to the peers a node will actually contact.

Records are Protocol Buffers, optionally persisted in bbolt with an
in-memory LRU in front~\cite{iptrackerbolt}. Kubernetes already speaks
protobuf; the distinction is the reduced schema and the topology. The
details are less flattering. The LRU admits updates as well as lookups,
so churn among peers a node never contacts can evict the ones it
contacts constantly; store access is serialized under a mutex, coupling
write bursts to lookup latency; the full sync clears local state and
reinserts the prefix, which is neither an atomic snapshot replacement
nor a constant-memory bootstrap; and a bounded object cache bounds
neither process nor machine memory once bbolt's mapping, page cache,
and full-sync buffers are counted. Issue 114659 documents why informer
memory is worth reducing~\cite{k8s114659}; it does not measure this
implementation.

The only quantitative data are single-client integration logs posted
with PR 218~\cite{iptracker218}: 10,000 records gave 643 ms average
propagation, 35 MB of server Go heap, and a 1 MB client store; 40,000
records at 30,000 writes/s gave 6.6 s and 174 MB, with the etcd
maintainer in the thread reporting watch stalls above that size; 10
million records throttled to 1,024 writes/s finished in 9,828 s at
43 ms average propagation, 24.7 GiB of server heap, and a 1.25 GiB
client store. One process, one machine, no packet path, an LRU sized to
hold every record: these bound neither agent memory under eviction nor
fan-out to hundreds of clients. They do show that the distribution
tier, not the client, holds the bulk of the state, and that in the
tested range propagation tracks write rate rather than record count.
Before persisted records serve as authorization metadata the design
owes answers on revision recovery after compaction, reflector failover,
disk loss, and above all IP reuse: a stale record for a reused address
is a policy bypass, not an availability delay. The IPTracker plugin
flavor also diverts all forwarded traffic to the queue
(\texttt{ManagedIPs} reports every address as managed), so its
$\lambda_q$ differs from the one in Section~\ref{sec:model}. IPTracker
is implemented future work; no result in Section~\ref{sec:evaluation}
used it or supports it.

\subsection{Testing the Architectural Hypothesis}

Section~\ref{sec:gaps} lists what is missing; this is the order in
which we would collect it, by how much each result could change the
conclusion. First, rebuild Kindnet with the agent metrics server
enabled and rerun one tier, preferably 7 identities at 5,000 Pods per
ReplicaSet; all six agent queries returned no samples, so no run has a
$\lambda_q$, a $\mu$, drop counts, or agent CPU and memory. Second,
collect the \texttt{wait-for-gateway} init container's
\texttt{Latency:} lines, which every sandbox already logs and the
harness did not keep; large values on the slow Pods implicate metadata
convergence or the verdict path, small ones exonerate both and point at
kubelet or the runtime. Third, equalize achieved Pod throughput or
explain why the \sys{} cluster created 103--124 Pods per second against
Cilium's 302--328 on identical control-plane settings; a slower creator
means a smaller per-node burst, which flatters the tail, so the
closeness at 700 and 3,500 identities proves less than it appears to.
Fourth, a matched pair at 35,000 identities, a Cilium raw-Pod mesh run
or a \sys{} run with one Pod per ReplicaSet. Fifth, repeat trials; with
one sample per tier the 0.5 s and 2.4 s differences are
indistinguishable from noise. Sixth, a \texttt{cilium-dbg bpf policy
get} dump and the effective \texttt{cilium-config} from a gateway node,
so that Figure~\ref{fig:scaling} rests on a measurement rather than the
indirect argument of Section~\ref{sec:evaluation}. Seventh, a
single-node connection microbenchmark with a denied-traffic mix, run
into overload in both queue modes, since the capacity boundary of
Section~\ref{sec:design} is an overload point nobody has
measured. Last, vary label sharing and workload freshness independently
at fixed live-Pod count, to separate identity dissemination from Pod
churn.

The disconfirming outcomes are easy to state. Agent saturation under
burst, or \texttt{Latency:} lines that account for the tail, would
attribute the 151 s and 421 s tails to deferred evaluation rather than
a confound; equalized throughput that widens the gap at 700 and 3,500
identities would remove the one place where \sys{} currently looks
competitive; a factored-kernel baseline completing the 35,000-identity
mesh cheaply would reduce the capacity claim to a statement about one
representation; and a denied-traffic microbenchmark that fills the
1,024-packet queue at a connection rate a real tenant could produce
would turn the overflow boundary from a constant in the source into a
measured limit.

\subsection{Limits of the Present Claim}

The design does not promise unlimited identities, constant total
memory, zero startup delay, or resource use independent of topology. It
removes one eager cross-product by keeping semantic evaluation in
userspace and caching the outcome per connection, a good trade only
where the admission and consistency costs are acceptable. The \sys{}
mesh sweep supports a capacity statement: a tier the Cilium harness
excluded on predicted per-endpoint map demand completes under \sys{}
with all 35,000 Pods Running and none stranded, and going from 700 to
3,500 identities does not raise its P99 (5.331 s to 2.791 s). It does
not support an admission-latency statement: \sys{} is slower than
Cilium at every matched tier, by 2.4 s and 0.5 s where it is close and
by a factor of 22 at 7 identities (151.479 s against 6.767 s); the
raw-Pod tier records 421.410 s with two startup-SLO failures; achieved
Pod throughput is 2.5--3$\times$ lower for reasons nobody recorded;
every tier is one trial; and the agent metrics are empty. Enforcement
is verified only on the positive path, by the gateway gate every
sandbox passes; the negative path was checked once per cluster by
\path{scripts/verify-netpol.sh}, unarchived.

Two further caveats. Results are version-specific on both sides: an
archived Cilium v1.18.6 run stranded 19 Pods at 700 identities on a
policy-cache race fixed in v1.20.0~\cite{cilium7515}, and the \sys{}
figures come from one Kindnet v1.0.1 build with no metrics port and
unrecorded agent flags. And the semantics enforced are looser than the
discussion sometimes implies: NetworkPolicy leaves the effect of a
policy update on established connections
implementation-defined~\cite{k8snetworkpolicy}, so caching an accepted
verdict for the life of the connection is conforming on its own, strict
mode's conntrack scan is a choice \sys{} makes rather than a
requirement the API imposes, and an engine that does revoke is equally
within the API. Comparing revocation cost compares design choices, not
conformance.

\section{Conclusion}
\label{sec:conclusion}

NetworkPolicy at scale is a joint problem of representation, placement,
and synchronization, and the three cannot be optimized one at a time. A
fast kernel lookup does not make workload admission cheap, and a flat
latency curve does not mean the policy state has no ceiling. NFQUEUE,
conntrack, and NRI permit a different decomposition: interpret the
policy when communication is attempted, keep the accepted decision in
the kernel's connection state, and learn local Pod addresses from the
runtime rather than from the API server. None of these mechanisms is
new, and the decomposition inherits the known weakness of reactive
processing: its cost concentrates in the miss bursts.

The archived experiments on 720 workers and 35,000 Pods show both
sides. Cilium's completed sweeps hold low post-scheduling latency and
flat worker CPU across their feasible region, halve achieved throughput
at 35,000 ReplicaSets, exclude the 35,000-identity bidirectional mesh
on predicted map capacity, and stranded 5,694 Pods in one hub-and-spoke
raw-Pod burst. The \sys{} sweep completes the excluded tier with every
Pod Running, stays within 2.4 s of Cilium's P99 at 700 and 3,500
identities, and records 151 s and 421 s tails in its two burst tiers
that the artifacts cannot attribute, at a Pod-creation rate
2.5--3$\times$ below Cilium's that they cannot explain either. Because
every sandbox gates on a permitted TCP connection to the gateway before
it reports Running, these latencies contain the first permitted
cross-node communication; that makes them a more honest measure than
container startup, and it makes the unexplained tails the most
important open item, since they are tails on the very quantity the
design claims to improve. Completing the comparison requires the
measurements of Table~\ref{tab:gaps}, beginning with an agent that
exposes its metrics, the init-container latencies the sandboxes already
record, and a second trial of the matched tiers.
